% Pursuing Stacks (cotes 134-2 to 134-8): archive cote 134-5, site page 10 (PDF page 11), typescript page 197.
% typescript. Draft compared: ps4.tex:425-486.
% AI-transcribed by Claude Opus 5.5 (claude-opus-5-5, adaptive thinking, effort high), Message Batches API, batch ps2 (msgbatch_01NVwr1Yuo8qJf5rykeGBXPK), 2026-11-12; not yet checked by eye.
\scanpage{134-5/10}
\mspage{197}
\inserted{to} the derived category of the abelian category of all \underline{A}-valued diagrams
of this type. In precise terms, for any small category I, we get
the categ\inserted{o}ry D(\underline{Hom}(I,\underline{A})), depending functorially in a contravariant
way on I. Rewriting this category $D_{\underline{A}}(I)$ say, the idea was to consider
(*) \[ I \mapsto D_{\underline{A}}(I) \quad , \]
possibly \inserted[typed]{with I suitably} restricted \xout{xx} (for instance to finite categories, or to finite
ordered sets, corresponding to finite \uline{commutative} diagrams), as embo-%
dying the "full" triangulated structure \unsure{defihed|defined} by \underline{A} . This of course
at once raises a number of questions, such as recovering the usual
triangulated structure of $D(\underline{A}) = D_{\underline{A}}(e)$ ( e the final object of (Cat))
in terms of (*), and pinning down too the relevant formal properties
(and possibly even extra structure) one had to assume on (*). I had
never so far taken the time to sit down and play around some and see
how this goes through, expecting that surely someone else would do it
some day and I would be informed - but apparently in the last eighteen
years nobody ever was interested. Also, it had been rather clear from
the start that Verdier's constructions could be adapted and \inserted{did} make sense
for non-commutative homotopy set-ups, which was also appearant in betw-%
ween the lines in Gabriel-Zisman's \xout{xxxxxxxxxf} book on the foundations
of homotopy theory, and a lot more explicitly in Quillen's axiomatizati-%
on of homotopical algebra. This \inserted{axiomatization} I found very appealing indeed - and
right now \inserted{still} his little book is my most congenial and main source of infor-%
mation on foundational matters of homotopical algebra. I reme\inserted{m}ber though
\inserted{my} being a little disappointed at Quillen's not caring either to pursue the
matter of what exactly a "non commutative triangulated category structu\xout{\unsure{re}}-%
re" (of the type he was getting from his model categories) was, just
contenting himself tomumble a few words about existence of "higher
structure" (t\over{hen}{han} just the Dold-Puppe sequences), which \inserted{(he implied)} \unsure{needed} to be under-%
stood. I felt of course that presumably the variance formalism (*) should
furnish any kind of "higher" structure one was looking for, but it was'
nt really my business to check.

It still is'nt, however I did some homework on (*) - it was the
first thing indeed I looked at in these six weeks, and some main featu-%
res came out very readily indeed. It turns out that the main formal
variance property to demand on (*), presumably even the only one , is
that for a given map \inserted{$f$:} I $\to$ J on the indexing categories of diagram-types
I and J , the corresponding functor
\[ \hand{f^*:} D_{\deleted{\underline{A}}}(J) \hand{\longrightarrow} D_{\deleted{\underline{A}}}(I) \]
\uline{should have both a left and a right adjoint}, \inserted{say $f_!$ and $f_*$.} In case J= e , the two
\margin{10}
